\documentclass[12pt]{article}
\usepackage[utf8]{inputenc}
\usepackage[T1]{fontenc}
\usepackage{hyperref}
\title{Flemish Sign Language Language Resources for K-12 Curriculum Development}
\author{Compiled by Claude Sonnet 5 (claude-sonnet-5) for Andr\'as Kornai}
\date{}
\begin{document}
\maketitle

\section{Introduction}
Flemish Sign Language (Vlaamse Gebarentaal, VGT; ISO 639-3 vgt) is the deaf community sign language of Flanders, Belgium. This report covers only what appeared in a small number of web searches. Little material was found, and no ready-made K-12 curriculum was found.

\section{Primary Educational Resources}
No VGT textbooks or primers for school use turned up in the searches. The Adviescommissie Vlaamse Gebarentaal (the advisory commission on VGT) reports that teachers and teaching materials for VGT initiatives in schools are often lacking. Doof Vlaanderen states that bilingual VGT--Dutch education within mainstream schooling does not currently exist in Flanders; interpreters are used only to give access to lessons in Dutch. Deaf children outside mainstream schooling attend deaf schools.
\begin{itemize}
\item Doof Vlaanderen and the advisory commission, joint press release: \url{https://www.doof.vlaanderen/sites/default/files/imce/persbericht_doof_vlaanderen_en_adviescommissie_vlaamse_gebarentaal_finaal.pdf}
\item Advisory commission on quality VGT teaching for hearing adults: \url{https://vlaanderen.be/cjm/sites/default/files/2023-01/Beleidsadviezen_m.b.t._een_kwalitatief_aanbod_onderwijs_VGT_voor_horende_volwassenen.pdf}
\item ACVGT advice of 14 November 2018: \url{https://vlaanderen.be/cjm/sites/default/files/2021-12/ACVGT-advies-2018.11.14.pdf}
\end{itemize}

\section{Linguistic and Reference Materials}
The Woordenboek Vlaamse Gebarentaal is an online dictionary of roughly 10,000 entries, searchable by Dutch keyword, handshape and location. The Global Signbank VGT dataset behind it is managed by the Vlaams Gebarentaal Centrum, and its licence is restricted. The advisory commission notes that little linguistic research on VGT exists and that few detailed grammars and dictionaries are available to teachers and students.
\begin{itemize}
\item Woordenboek Vlaamse Gebarentaal (entry in the sign-lang compendium): \url{https://www.sign-lang.uni-hamburg.de/lr/compendium/lex/vgtdictionary.html}
\item Global Signbank VGT: \url{https://www.sign-lang.uni-hamburg.de/lr/compendium/lex/globalsignbankvgt.html}
\item sign-lang@LREC Anthology entry on the VGT dictionary: \url{https://www.sign-lang.uni-hamburg.de/lrec/data/vgtdictionary.html}
\item Abstract on lexicographic projects and VGT (KU Leuven): \url{https://lirias.kuleuven.be/bitstream/123456789/522094/1/Abstract+Lexicographic+projects+and+VGT.pdf}
\end{itemize}

\section{Religious and Community Resources}
Doof Vlaanderen, the deaf community organisation, appears above. No religious resources were found in the searches.

\section{Digital and Audio Resources}
The online dictionary and Global Signbank are the only digital resources found. No video lesson series for school use was found.

\section{Curriculum Development Context}
The sources describe a gap. There is no bilingual VGT--Dutch mainstream education, teaching materials are scarce, and little research has been turned into classroom materials. A VGT-medium K-12 curriculum would largely have to be built from scratch, working with the deaf community and the advisory commission.
\end{document}